\documentclass[11pt,a4paper]{article}
\usepackage[T1]{fontenc}
\usepackage[utf8]{inputenc}
\usepackage{lmodern,amsmath,amssymb}
\usepackage[margin=25mm]{geometry}
\usepackage{longtable,booktabs,array,calc}
\usepackage[hidelinks]{hyperref}
\hypersetup{pdftitle={SU(3) in four dimensions: every constant of this programme, evaluated},pdfauthor={navstokgap research working notes}}
\newcommand{\tightlist}{\setlength{\itemsep}{0pt}\setlength{\parskip}{0pt}}
\setlength{\emergencystretch}{3em}
\setcounter{secnumdepth}{0}
\begin{document}
\hypertarget{su3-in-four-dimensions-every-constant-of-this-programme-evaluated}{%
\section{SU(3) in four dimensions: every constant of this programme,
evaluated}\label{su3-in-four-dimensions-every-constant-of-this-programme-evaluated}}

With the goal narrowed to \(SU(3)\), each bound proved in these notes
becomes a number. The group data are \[N=3,\quad \dim G=8,\quad
C_2(R_{\rm f})=\frac{N^2-1}{2N}=\frac43,\quad
C_2(\text{adj})=N=3,\quad
Z(G)=\mathbb Z_3,\] in the normalization
\(\operatorname{tr}T^aT^b=\tfrac12\delta^{ab}\), and the one- and
two-loop coefficients are
\[b_0=\frac{11N}{48\pi^2}=\frac{11}{16\pi^2}=0.06966,\qquad
\frac{b_1}{2b_0^2}=\frac{51}{121}=0.4215 .\] The consequences,
collected: the strong-coupling criterion of
\href{strong-coupling-uniform-gap.md}{the T2 note} reads
\(54/g^4\le\beta_*\); the gap it gives is
\(\Delta_{a,L}\ge\tfrac23\gamma g^2\,\hbar c/a\); the unperturbed
flux-loop gap is \(\tfrac83g^2\hbar c/a\); the Lieb--Robinson velocity
is \(v\le96e\,c/g^2\simeq261\,c/g^2\); the exact identities are
\(\Delta V=\tfrac{16}3(3|\mathcal P|-V)\) and \(|\nabla V|^2\le48V\);
the number of renormalization doublings from a bare coupling
\(g_{\rm UV}^2=1/2\) to a strong-coupling threshold is
\[n\simeq\frac{1}{2b_0\log2}\Big(\frac1{g_{\rm UV}^2}-\frac1{g_{\rm thr}^2}\Big)\le\frac{2}{0.0966}\simeq21,\]
fewer than the thirty of \(SU(2)\) because the coupling runs faster; and
the flowed upper bound \(4.26\,\hbar c/\sqrt{8t}\) is unchanged, because
in the free limit the eight gluon species multiply numerator and
denominator equally. Two features of \(SU(3)\) that \(SU(2)\) lacks are
recorded in Section 4: the defining representation is complex, so
\(\operatorname{tr}U_p\) is complex and the plaquette term keeps only
its real part, and the centre is \(\mathbb Z_3\), so a periodic box
carries \(27\) electric-flux sectors instead of \(8\). Constants
explicit; nothing promoted.

\hypertarget{group-data}{%
\subsection{1. Group data}\label{group-data}}

\begin{longtable}[]{@{}
  >{\raggedright\arraybackslash}p{(\columnwidth - 4\tabcolsep) * \real{0.3333}}
  >{\raggedright\arraybackslash}p{(\columnwidth - 4\tabcolsep) * \real{0.3333}}
  >{\raggedright\arraybackslash}p{(\columnwidth - 4\tabcolsep) * \real{0.3333}}@{}}
\toprule\noalign{}
\begin{minipage}[b]{\linewidth}\raggedright
quantity
\end{minipage} & \begin{minipage}[b]{\linewidth}\raggedright
value
\end{minipage} & \begin{minipage}[b]{\linewidth}\raggedright
where used
\end{minipage} \\
\midrule\noalign{}
\endhead
\bottomrule\noalign{}
\endlastfoot
\(N\) & \(3\) & everywhere \\
\(\dim G\) & \(8\) & free-limit species count \\
\(C_2(R_{\rm f})\) & \(4/3\) & flux-loop energy, T2 gap \\
\(C_2(\text{adj})\) & \(3\) & Jacobian curvature factor \\
centre & \(\mathbb Z_3\) & flux sectors, valley minima \\
\(\operatorname{tr}(1)\) in \(R_{\rm f}\) & \(3\) & plaquette norm
\(\|w_p\|\le4N\hbar c/(ag^2)=12\hbar c/(ag^2)\) \\
\(b_0\) & \(11/(16\pi^2)=0.06966\) & running, step count \\
\(b_1/(2b_0^2)\) & \(51/121=0.4215\) & the power in
\(a\Lambda_{\rm lat}\) \\
\end{longtable}

The lattice \(\Lambda\) parameter of
\href{mass-gap-obligations-lattice.md}{the obligations map} §5 is
therefore
\[a\Lambda_{\rm lat}(g)=\big(b_0g^2\big)^{-51/121}\,e^{-1/(2b_0g^2)}\big[1+O(g^2)\big],
\qquad \frac1{2b_0}=\frac{8\pi^2}{11}=7.180 .\]

\hypertarget{the-proved-bounds-evaluated}{%
\subsection{2. The proved bounds,
evaluated}\label{the-proved-bounds-evaluated}}

\textbf{T1.} Unchanged in form: unique positive physical ground state
and \(\delta(g;N_s,SU(3))>0\) at every finite lattice and coupling.

\textbf{T2} (\href{strong-coupling-uniform-gap.md}{T2 note}). The
perturbation parameter is
\[\beta=\frac{48N^2}{g^4(N^2-1)}=\frac{48\cdot9}{8\,g^4}=\frac{54}{g^4},\]
so the criterion \(\beta\le\beta_*(3,\{0,1\}^3)\) becomes
\[g\ \ge\ g_0=\Big(\frac{54}{\beta_*}\Big)^{1/4},\] and the gap is
\[\Delta_{a,L}\ \ge\ \gamma\,\frac{g^2}{2}\,\frac43\,\frac{\hbar c}{a}
=\frac{2\gamma}{3}\,g^2\,\frac{\hbar c}{a}.\] The unperturbed physical
excitation, a single plaquette flux loop in the fundamental, costs
\(2C_2g^2\hbar c/a=\tfrac83g^2\hbar c/a\).

\textbf{Feynman--Bijl, strong coupling}
(\href{lattice-gap-upper-bounds.md}{upper-bound note} Corollary 3). At
\(g=\infty\) the Haar variance is
\(\operatorname{Var}(\operatorname{Re}\operatorname{tr}U_p)=\tfrac12\)
for \(N\ge3\), so the two-sided statement reads
\[\frac{2\gamma}{3}g^2\,\frac{\hbar c}{a}\ \le\ \Delta^{\rm phys}_{a,L}\ \le\ 12\,g^2\,\frac{\hbar c}{a}\]
for \(g\) beyond both thresholds, a ratio of \(18/\gamma\) between the
two sides.

\textbf{Lieb--Robinson} (\href{lieb-robinson-kogut-susskind.md}{note}).
Four plaquettes per link and \(\|w_p\|\le12\hbar c/(ag^2)\) give
\(J\le48\hbar c/(ag^2)\) and
\[v=\frac{2eaJ}{\hbar}\ \le\ \frac{96\,e\,c}{g^2}\ \simeq\ \frac{261\,c}{g^2},\]
so the lattice light cone exceeds \(c\) for \(g^2<261\), that is
throughout any regime of interest, and the strong-coupling correlation
length is \(\xi\le C'\,a\,/(\gamma g^4)\cdot\tfrac94\).

\textbf{Exact identities} (\href{magnetic-energy-identities.md}{note}).
With \(C_2=4/3\) and \(N=3\),
\[\Delta V=\frac{16}{3}\big(3|\mathcal P|-V\big),\qquad
|\nabla V|^2\le16NV=48\,V,\] and the ground-state sum rule of that note
holds with \(2C_2=8/3\).

\textbf{Blocking} (\href{blocking-criterion-monotone.md}{criterion
note}). The block criterion is
\(\beta_{\rm block}=72M^2/(g^2\delta(g;M))\) and the direct one is
\(\beta_{\rm direct}=24/g^4\), so the ratio at strong coupling is
\(3M^2/8\) as before, blocking losing by \(3/2\) already at \(M=2\).

\textbf{Small volume} (\href{low-dimensional-mass-gap.md}{G07}
Proposition 7). The constant-mode sector has \(D=3\) and gauge group
\(SU(3)\), so the matrix model runs over \(\vec x_i\in\mathbb R^8\) with
the structure constants of \(SU(3)\) in place of the cross product; the
gap keeps the form \(\delta_1^{(3)}g^{2/3}\hbar c/L\) with a pure number
\(\delta_1^{(3)}\) specific to the group.

\hypertarget{the-step-count}{%
\subsection{3. The step count}\label{the-step-count}}

With \(2b_0\log2=2(0.06966)(0.6931)=0.09657\), the number of
renormalization doublings needed to run from a bare coupling to a
strong-coupling threshold is
\[n=\frac{1}{0.09657}\Big(\frac1{g_{\rm UV}^2}-\frac1{g_{\rm thr}^2}\Big).\]

\begin{longtable}[]{@{}lll@{}}
\toprule\noalign{}
\(g_{\rm UV}^2\) & \(n\) to \(1/g_{\rm thr}^2\to0\) & scale ratio
\(2^n\) \\
\midrule\noalign{}
\endhead
\bottomrule\noalign{}
\endlastfoot
\(1\) & \(10.4\) & \(1.4\times10^{3}\) \\
\(1/2\) & \(20.7\) & \(1.7\times10^{6}\) \\
\(1/4\) & \(41.4\) & \(2.9\times10^{12}\) \\
\end{longtable}

For \(SU(2)\) the same table reads \(15.5\), \(31.1\), \(62.1\), since
\(b_0\) is smaller by \(2/3\). \textbf{The \(SU(3)\) problem needs about
a third fewer steps than \(SU(2)\)} at the same bare coupling, which is
the only quantitative respect in which the physical case is easier.

\hypertarget{two-features-specific-to-su3}{%
\subsection{\texorpdfstring{4. Two features specific to
\(SU(3)\)}{4. Two features specific to SU(3)}}\label{two-features-specific-to-su3}}

\textbf{The defining representation is complex.} For \(SU(2)\),
\(\operatorname{tr}U_p\) is real, and the plaquette term is the full
trace; for \(SU(3)\) it is complex and the Wilson term keeps
\(\operatorname{Re}\operatorname{tr}U_p\). Consequences: the bound
\(\int(\operatorname{tr}U)^2dU=0\) used in
\href{lattice-gap-upper-bounds.md}{the upper-bound note} Corollary 3
holds for \(N\ge3\) and gives Haar variance \(\tfrac12\) rather than
\(1\); and charge conjugation is a nontrivial symmetry, so the spectrum
splits into \(C\)-even and \(C\)-odd sectors, and the glueball channel
relevant to the gap is \(C\)-even.

\textbf{The centre is \(\mathbb Z_3\).} The valley minima of
\href{torus-valley-potential.md}{the valley note} §4 lie at holonomies
that are central, so there are three per direction and \(27\)
electric-flux sectors on a three-torus, against \(8\) for \(SU(2)\). The
tunnelling barrier between them is the same \(O(\hbar c/L)\), and the
sector structure is what the small-volume spectrum organizes itself by.

\hypertarget{what-does-not-change}{%
\subsection{5. What does not change}\label{what-does-not-change}}

The structural results are group-independent and stand as written:

\begin{itemize}
\tightlist
\item
  the seven closed routes of \href{mass-gap-position.md}{the position
  note} §3, each closed by an argument that uses no property of the
  group beyond compactness and non-abelianness;
\item
  the moment hierarchy and the reduction of \(m<\infty\) to one flowed
  correlator, where the free-limit value \(4.26\,\hbar c/\sqrt{8t}\) is
  unchanged because the \(\dim G=8\) species multiply the susceptibility
  and the gradient term equally;
\item
  the equivalence of T2\('\) with the absence of a zero-temperature
  phase transition;
\item
  the three places where the non-abelian structure is isolated, all of
  which are statements about the commutator and hold for \(SU(3)\) with
  the \(SU(3)\) structure constants.
\end{itemize}

\hypertarget{consequence-for-state}{%
\subsection{6. Consequence for STATE}\label{consequence-for-state}}

Specializing to \(SU(3)\) fixes every constant and shortens the step
count to about twenty-one doublings from \(g_{\rm UV}^2=1/2\). It
changes no structural conclusion, and the open obligations are those of
\href{mass-gap-position.md}{the position note} §5. The next work should
be carried out with these numbers in place, and any new bound should be
reported in them.
\end{document}
